\section{Introduction}
\label{sec:intro}
Prediction-powered inference (PPI) turns a small labelled sample and a large unlabelled one, scored by a fixed machine-learning predictor $f$, into a valid confidence interval that is shorter than the labelled-only interval
\citep{angelopoulos2023ppi}; PPI++ tunes the weight of the predictor and is, among corrections that are affine in $f$, the shortest of them \citep{angelopoulos2023ppiplus}.
We ask a question that sits after that: \emph{with the labelled set, the unlabelled set and the predictor all held fixed, and nominal coverage required, can the PPI/PPI++ interval be made shorter, by how much, and under what conditions?}
Nothing is bought here with more labels, a better $f$ or a different labelling design; the only freedom is what the analyst does with the same records.

The answer we find is conditional, and for the most common use it is negative. Writing the estimator as a control variate (Sec.~\ref{sec:theory}), PPI++ is the best affine function of $f$, and a known efficiency bound limits any control variate built from $(f,X)$; the headroom is the gap between the two.
On the seven real pools of the \texttt{ppi-python} package that gap, measured for the adjuster classes we tried (fitted on half-pools, so a lower bound on the true gap), is at most about 1.2\% of the interval width for the mean (Sec.~\ref{sec:headroom}), and learning a better control variate from the $n$ labels costs more than it returns:
of 270 (pool, $n$, adjuster, calibration) cells on the four main pools at $n\le1000$, 252 are significantly \emph{wider} than PPI++ calibrated the same way and none is more than 0.55\% narrower (Sec.~\ref{sec:real}).
Where the headroom is large, which we can only establish in populations with known truth, it is realised, but only above a threshold $n$ and only by an adjuster that matches its shape: up to 26\% narrower when $E[Y|f]$ is far from affine, and
22\% to 29\% when features $X$ carry signal that $f$ misses, for $n$ from 50 (a linear adjuster on $(f,X)$) to several hundred (flexible adjusters), at coverage within 1.5 points of PPI++'s (Sec.~\ref{sec:synth}). On a regression problem with real data (OLS of income on covariates) the saving is 2.8\% at $n=1000$ (Sec.~\ref{sec:ols}).

\paragraph{Contributions.}
(i) A variance calculus that places PPI++ inside a family of cross-fitted control variates, states the efficiency floor for the mean (a known bound, restated) and a first-order account of the estimation cost (Sec.~\ref{sec:theory}).
(ii) A like-for-like protocol: every width is a paired ratio against PPI++ with the same interval calibration (Wald, $t$ or studentised bootstrap) and is read together with coverage and its Wilson interval, because a narrower interval that undercovers is not a saving (Secs.~\ref{sec:method}, \ref{sec:validity}).
(iii) The measured headroom, finite-sample cost and break-even $n$ on the real pools, in synthetic populations whose theory is known, and for OLS coefficients (Secs.~\ref{sec:headroom}--\ref{sec:ols}).
(iv) A statement of when the answer is ``no'' (calibrated or near-affine $f$, no informative $X$, small $n$, label shift) and every cell behind it, with the code that reproduces them.
Positive and negative results are reported with equal weight; the claims are limited to the estimands, adjuster classes and sampling protocol described in Secs.~\ref{sec:method}--\ref{sec:exp} and in the discussion.
